\documentclass{article}
\usepackage{amsmath, amssymb}

\title{1.2.3-35}

\begin{document}
\maketitle

(a) $\sup_{R(i)}a_i + \sup_{S(j)}b_j = \sup_{R(i)}\sup_{S(j)}(a_i + b_j)$

If $a_i, b_j \geq 0: \sup_{R(i)}a_i \cdot \sup_{S(j)}b_j = \sup_{R(i)}\sup_{S(j)}(a_ib_j)$

(b) $\sup_{R(i)} a_i = \sup_{R(j)} a_j = \sup_{R(p(j))} a_{p(j)}$

(c) $\sup_{R(i)}\sup_{S(j)} a_{ij} = \sup_{S(j)} \sup_{R(i)} a_{ij}$

(d) $\max(\sup_{R(j)} a_j, \sup_{S(j)} a_j ) = \sup_{R(j) \lor S(j)} a_j$

If $S(j)$ never satisfies: $\max(\sup_{R(j)} a_j, \sup_{S(j)} a_j ) = \sup_{R(j)} a_j$

$\implies \sup_{S(j)} a_j \leq \sup_{R(j)} a_j, \forall R \implies \sup_{S(j)} a_j = -\infty$


\end{document}
